\section{第7章\quad EM 算法}

\subsection{故事：缺一块拼图的极大似然}

极大似然估计（MLE）的流程很顺：写似然 $\to$ 求导 $\to$ 置零解出 $\hat\theta$。
可一旦数据\textbf{缺了一块}——有的观测丢失了、或者每个数据"来自哪个子总体"
这个标签看不见——似然函数里就会出现"对缺失部分求和/积分"的项，
变成一团解不开的对数和（$\log\sum$），求导置零没有解析解。

EM 算法（Expectation--Maximization, Dempster et al.\ 1977）的破局思路像
"鸡生蛋、蛋生鸡"的迭代：
\begin{center}
\fbox{\parbox{0.9\textwidth}{\centering
\textbf{E 步}：假装参数已知，把看不见的缺失信息"猜"出来（求条件期望）；\\
\textbf{M 步}：假装缺失信息已知，像普通 MLE 一样更新参数（求最大化）。\\
两步轮流做，似然只升不降，直到收敛。}}
\end{center}

\subsection{一般框架}

\paragraph{记号。}观测数据 $X$（不完全数据），缺失/隐变量 $Z$，
完全数据 $(X,Z)$ 的对数似然 $\log L(\theta;X,Z)$ 通常\textbf{很好优化}。

\begin{kbox}{EM 算法（必背）}
给定初值 $\theta^{(0)}$，第 $i$ 次迭代：
\begin{itemize}[leftmargin=2em]
  \item \textbf{E 步}：算 Q 函数——完全数据对数似然在当前参数下对缺失变量的条件期望
  \[
  Q\big(\theta,\theta^{(i-1)}\big)
  =\E_{Z\mid X,\theta^{(i-1)}}\big[\log L(\theta;X,Z)\big];
  \]
  \item \textbf{M 步}：$\theta^{(i)}=\arg\max_\theta Q\big(\theta,\theta^{(i-1)}\big)$。
\end{itemize}
重复至 $\|\theta^{(i)}-\theta^{(i-1)}\|$ 或似然增量小于阈值。
\end{kbox}

\begin{pbox}
Q 函数看着吓人，其实就一句话：\textbf{缺失的部分不知道，就用"当前参数下的最佳猜测"
代进去，把完全数据的对数似然变成一个能算的普通函数}。
类比填空题：考卷（似然）有几个空（隐变量）看不清，先按目前的理解把空填上（E 步），
再按填好的卷子认真作答（M 步）；答完理解加深了，回头重填空、重作答，如此循环，
每轮都比上轮更好（单调性）。
\end{pbox}

\begin{kbox}{两条性质（概念题常考）}
\begin{enumerate}[leftmargin=2em]
  \item \textbf{单调性}：每次迭代观测数据似然不降，
  $L(\theta^{(i)};X)\ge L(\theta^{(i-1)};X)$，故必收敛（似然有上界时）；
  \item \textbf{局部性}：只保证收敛到\textbf{驻点/局部最大}，不保证全局最优
  $\Rightarrow$ 实践中要\textbf{多个初值}多跑几次取最优。
\end{enumerate}
\end{kbox}

\paragraph{考试套路。}EM 大题基本都是同一个剧本：
(1) 写出完全数据似然；(2) E 步算出隐变量的条件期望/后验概率；
(3) M 步求导得参数更新公式；(4) 也许让你代一步数值。下面用四个经典模型走一遍。

\subsection{例一：两枚硬币（隐标签问题）}

\begin{exbox}[两硬币模型]
有 A、B 两枚硬币，正面概率 $p_A,p_B$ 未知。做 $n$ 轮实验：每轮随机摸一枚
（不知道摸的是谁），抛 $m$ 次记录正面数 $x_i$。写出 EM 算法。
\end{exbox}
\begin{sol}
隐变量 $Z_i\in\{A,B\}$：第 $i$ 轮用的是哪枚（设摸到各为 $1/2$）。

\textbf{E 步}（贝叶斯公式算"软标签"）：当前参数 $p_A^{(t)},p_B^{(t)}$ 下，
第 $i$ 轮来自 A 的后验概率
\[
\mu_i=\frac{\binom{m}{x_i}(p_A^{(t)})^{x_i}(1-p_A^{(t)})^{m-x_i}}
{\binom{m}{x_i}(p_A^{(t)})^{x_i}(1-p_A^{(t)})^{m-x_i}+\binom{m}{x_i}(p_B^{(t)})^{x_i}(1-p_B^{(t)})^{m-x_i}} .
\]

\textbf{M 步}（按软标签加权的"正面频率"）：
\[
p_A^{(t+1)}=\frac{\sum_i\mu_i x_i}{\sum_i\mu_i\,m},\qquad
p_B^{(t+1)}=\frac{\sum_i(1-\mu_i)x_i}{\sum_i(1-\mu_i)\,m}.
\]
直觉：若知道标签，$\hat p_A$ 就是 A 轮的正面比例；不知道，就按后验概率"打折计票"。
\end{sol}

\begin{pbox}
"软标签"$\mu_i$ 怎么理解？它不是硬说"第 $i$ 轮就是 A 硬币"，而是说
"按目前的参数看，这轮有 $\mu_i$ 的把握是 A"。于是 M 步计票时，
这轮的 $x_i$ 个正面不是全给 A 也不是全给 B，而是按 $\mu_i:(1-\mu_i)$ 分给两边——
就像一笔账说不清是谁的，就按可能性比例分摊。
\end{pbox}

\subsection{例二：多项分布的 197 只动物（经典必考）}

\begin{exbox}[Rao 的遗传学数据]
197 只动物分四类，频数 $\bm y=(y_1,y_2,y_3,y_4)=(125,18,20,34)$，
概率为
\[
\Big(\frac12+\frac\theta4,\ \frac{1-\theta}4,\ \frac{1-\theta}4,\ \frac\theta4\Big).
\]
直接 MLE 需解二次方程；用 EM 求 $\theta$。
\end{exbox}
\begin{sol}
\textbf{拆分技巧}：第一类概率 $\frac12+\frac\theta4$ 是"和"，把它拆成两个子类——
$z_1$（概率 $\frac12$，与 $\theta$ 无关）和 $z_2$（概率 $\frac\theta4$），
$z_1+z_2=y_1=125$。$z_2$ 即隐变量。

\textbf{完全数据似然}：$L\propto\theta^{z_2+y_4}(1-\theta)^{y_2+y_3}$，
其 MLE 是漂亮的比例形式 $\hat\theta=\frac{z_2+y_4}{z_2+y_4+y_2+y_3}$。

\textbf{E 步}：给定 $\theta^{(t)}$，$z_2\mid y_1\sim B\big(125,\ \frac{\theta^{(t)}/4}{1/2+\theta^{(t)}/4}\big)$，
\[
z_2^{(t)}=\E[z_2\mid y_1]=125\cdot\frac{\theta^{(t)}/4}{1/2+\theta^{(t)}/4}
=\frac{125\,\theta^{(t)}}{2+\theta^{(t)}}.
\]

\textbf{M 步}：
\[
\theta^{(t+1)}=\frac{z_2^{(t)}+34}{z_2^{(t)}+34+18+20}=\frac{z_2^{(t)}+34}{z_2^{(t)}+72}.
\]

\textbf{数值演示}：取 $\theta^{(0)}=0.5$：
$z_2^{(0)}=\frac{62.5}{2.5}=25$，$\theta^{(1)}=\frac{59}{97}\approx0.608$；
再迭代 $z_2^{(1)}=\frac{125\times0.608}{2.608}\approx29.15$，
$\theta^{(2)}\approx\frac{63.15}{101.15}\approx0.624$；……
收敛到 $\hat\theta\approx0.627$。
\end{sol}

\subsection{例三：正态样本含缺失值}

\begin{exbox}[缺失数据下估正态参数]
$X_1,\dots,X_n\overset{iid}{\sim}N(\mu,\sigma^2)$，但只观测到前 $m$ 个，
后 $n-m$ 个缺失。写出 EM 更新公式。
\end{exbox}
\begin{sol}
完全数据 MLE 需要两个充分统计量：$\sum_{i=1}^n x_i$ 和 $\sum_{i=1}^n x_i^2$。

\textbf{E 步}：对缺失的 $x_j$，在当前参数下
\[
\E[x_j\mid\cdot]=\mu^{(t)},\qquad
\E[x_j^2\mid\cdot]=\big(\mu^{(t)}\big)^2+\big(\sigma^{(t)}\big)^2 .
\]

\textbf{M 步}：
\[
\mu^{(t+1)}=\frac{1}{n}\Big(\sum_{i=1}^m x_i+(n-m)\mu^{(t)}\Big),
\]
\[
\big(\sigma^{(t+1)}\big)^2=\frac{1}{n}\Big(\sum_{i=1}^m x_i^2
+(n-m)\big[(\mu^{(t)})^2+(\sigma^{(t)})^2\big]\Big)-\big(\mu^{(t+1)}\big)^2 .
\]
\end{sol}

\begin{wbox}{头号易错点}
E 步填补的是\textbf{充分统计量的期望}，不是简单"把缺失值填成 $\mu^{(t)}$ 完事"！
$\E[x_j^2]=\mu^2+\sigma^2\ne\mu^2$——漏掉 $\sigma^2$ 项会让方差估计系统性偏小。
（这正是"填补均值再当完整数据算"这种朴素做法的错误所在。）
\end{wbox}

\subsection{例四：高斯混合模型（GMM）}

\paragraph{背景故事。}一群人的身高直方图呈双峰——其实是男女两个正态混在一起：
\[
f(x)=\sum_{k=1}^K\pi_k\,\varphi\big(x;\mu_k,\sigma_k^2\big),\qquad
\sum_k\pi_k=1 .
\]
隐变量 $Z_i$：第 $i$ 个人属于哪个成分。这是两硬币问题的连续版，EM 的"代表作"。

\begin{kbox}{GMM 的 EM（必背）}
\textbf{E 步}：算\textbf{责任}（responsibility）——第 $i$ 个观测属于第 $k$ 成分的后验概率
\[
\gamma_{ik}=\frac{\pi_k\,\varphi(x_i;\mu_k,\sigma_k^2)}
{\sum_{l=1}^K\pi_l\,\varphi(x_i;\mu_l,\sigma_l^2)} .
\]
\textbf{M 步}：按责任加权更新（记 $N_k=\sum_i\gamma_{ik}$，第 $k$ 成分的"有效人数"）
\[
\pi_k=\frac{N_k}{n},\qquad
\mu_k=\frac{1}{N_k}\sum_i\gamma_{ik}x_i,\qquad
\sigma_k^2=\frac{1}{N_k}\sum_i\gamma_{ik}(x_i-\mu_k)^2 .
\]
形式上就是"加权频率、加权均值、加权方差"——非常好记。
\end{kbox}

\begin{exbox}[GMM 单步手算]
两成分 GMM，当前参数 $\pi_1=\pi_2=0.5$，$\mu_1=0,\mu_2=4$，$\sigma_1=\sigma_2=1$。
对观测 $x=1$ 算责任 $\gamma_{11},\gamma_{12}$。
\end{exbox}
\begin{sol}
\[
\gamma_{11}=\frac{0.5\,\varphi(1;0,1)}{0.5\,\varphi(1;0,1)+0.5\,\varphi(1;4,1)}
=\frac{e^{-1/2}}{e^{-1/2}+e^{-9/2}}
=\frac{1}{1+e^{-4}}\approx0.982 ,
\]
$\gamma_{12}\approx0.018$。$x=1$ 离 $\mu_1=0$ 近得多，几乎全部"责任"归第 1 成分。
M 步时 $x=1$ 将以权重 0.982 参与 $\mu_1$ 的更新、以 0.018 参与 $\mu_2$。
\end{sol}

\subsection{本章考点地图}

\begin{ebox}{高频题型清单}
\begin{enumerate}[leftmargin=2em]
  \item \textbf{默写 EM 框架}：Q 函数定义、E/M 两步、单调性与局部最优（概念+简答）；
  \item \textbf{197 动物题}：拆分隐变量、推 E/M 公式、迭代一两步数值（最经典大题）；
  \item \textbf{两硬币/GMM}：后验概率（责任）公式+加权更新公式；
  \item \textbf{缺失正态数据}：E 步填充分统计量（二阶矩别漏 $\sigma^2$！）；
  \item 概念题：EM 为什么不会下降？为什么要多初值？E 步"软分配"与 K-means"硬分配"的关系。
\end{enumerate}
\end{ebox}
